
The cost is a lower wage within the son's occupation in mid-career; it is not where he starts and not his own
instability. Splitting his log wage into the CPS wage score of his occupation and his wage within it, the layoff's change
of $-0.603^{***}$ between 25--29 and 45--49 is $-0.392^{**}$ within the occupation and $-0.171^{**}$ between occupations
(panel C of Table~\ref{tab:layoff_mechanism}). Appendix~\ref{sec:skill} finds that the job a son holds at 25--34 predicts
how far he pulls away; the layoff does not work through it. A teen layoff lowers the chance of starting in a well-paid,
fast-growing job by only 0.041 ($p = 0.15$), and holding the son's early industry and occupation group fixed leaves the
within-son cost at $-0.225^{**}$ (against $-0.240^{**}$). Nor does the son repeat his father's instability: his own
displacement, unemployment, job changes and same-employer wage growth do not move with his father's teen layoffs. The one
trace is tenure, shorter by 4.2 years per unit share at 45--49 ($p = 0.093$), about seven months for one layoff
year.\footnote{Code: \texttt{106.LayoffCareerMechanism.py}.} Schooling, whose return rises with age
(Section~\ref{sec:ageprofile}), carries a sixth of the cost: a teen layoff barely changes a son's schooling ($-0.19$
years, s.e.\ 0.37), and holding schooling and its rising return fixed moves the within-son slope from $-0.236^{**}$ to
$-0.190^{**}$ and the change between 25--29 and 45--49 from $-0.914^{***}$ to $-0.763^{***}$.\footnote{Code:
\texttt{115.LayoffSchooling.py}.} What lowers the pay of these sons within their
occupation in their forties---skills, effort, health, or the jobs within an occupation they hold---the data here do not
separate.


